% !PRINTED 219
\weberpage{219}
Hieraus folgt, daß die Einteilung in Reihen gänzlich unabhängig
ist von der Willkürlichkeit in der Annahme über die
Wurzeln $\mathrm{sn}\,\Omega_1, \mathrm{sn}\,\Omega_2 \ldots$

\gesperrt{Es folgt aber noch weiter daraus, daß durch die
Substitutionen der Gruppe $\mathfrak{A}$ die Reihen nicht auseinandergerissen,
sondern nur untereinander vertauscht werden.}
Die Gruppe der Teilungsgleichung ist also imprimitiv, und die
einzelnen Reihen sind die Systeme der Imprimitivität (Bd.\,I, \S\,158).

Denn wendet man auf $(\mu,\mu')$ und $(\nu,\nu')$ gleichzeitig die
Substitution
\[
\begin{pmatrix} a, & b\\ c, & \partial \end{pmatrix}
\]
an, so bleibt die Kongruenz (3) erhalten.

Sucht man unter den Substitutionen der Gruppe $\mathfrak{A}$ die Substitutionen
aus, welche die Wurzeln einer Reihe nur untereinander
vertauschen, so erhält man eine Gruppe, und zwar einen Teiler
von $\mathfrak{A}$. Zu jeder Reihe gehört also ein solcher Teiler von $\mathfrak{A}$.
Bezeichnen wir mit $R_{\mu,\mu'}$ die Reihe, in der die Wurzel $x_{\mu,\mu'}$ vorkommt,
und mit $\mathfrak{A}_{\mu,\mu'}$ den zu dieser Reihe gehörigen Divisor von $\mathfrak{A}$,
so sind nach (3) [vgl. \S\,63, (3)] die Substitutionen
\[
\begin{pmatrix} a, & b\\ c, & \partial \end{pmatrix}
\]
von $\mathfrak{A}_{\mu,\mu'}$ durch die Kongruenz

\weq{4}{(\partial\mu - b\mu')\mu' + (c\mu - a\mu')\mu \equiv 0 \pmod{n}}
charakterisiert.

Ebenso erhält man eine Gruppe $\mathfrak{B}_{\mu,\mu'}$, wenn man die Substitutionen
in $\mathfrak{B}$ aufsucht, die die Wurzeln der Reihe $R_{\mu,\mu'}$ nur
unter sich vertauschen, deren Substitutionen:
\[
\begin{pmatrix} \alpha, & \beta\\ \gamma, & \delta \end{pmatrix}
\]
durch die beiden Kongruenzen:

\weq{5}{\begin{aligned}
(\delta\mu - \beta\mu')\mu' + (\gamma\mu - \alpha\mu')\mu &\equiv 0\\
\alpha\delta - \beta\gamma &\equiv 1
\end{aligned}\ \ (\mathrm{mod}\ n)}
charakterisiert sind.

Beispielsweise bestehen die Gruppen $\mathfrak{A}_{1,0}$, $\mathfrak{B}_{1,0}$ aus den Substitutionen
\[
\begin{pmatrix} a, & b\\ 0, & \partial \end{pmatrix}, \quad
\begin{pmatrix} \alpha, & \beta\\ 0, & \delta \end{pmatrix}, \quad \alpha\delta \equiv 1 \pmod{n}.
\]
